% =============================================================================
%  CHAPTER 22 -- PHASE 9: NONLINEAR EXTENSION + RESET
% =============================================================================
\section{Phase 9 --- Nonlinear signal extraction and the RESET test}
\label{ch:phase9}

\textbf{Reproduce it:} \code{python3 scripts/test\_nonlinear.py}, then \code{python3
scripts/analysis/compare\_nonlinear.py \&\& python3
scripts/plotting/render\_nonlinear.py}.\\
\textbf{Math used:} all of Chapter~\ref{ch:nonlinear}.\\
\textbf{Artifacts:} \code{scripts/model/nonlinear\_models.py},
\code{scripts/test\_nonlinear.py},
\code{results/tables/nonlinear\_\{summary,reset,learning,trading\}.csv},
\code{results/figures/nonlinear\_\{ic,reset,sharpe,learning\}.png},
\code{report/tables/nonlinear.tex}.

\subsection{Goal}

Answer one question honestly: \emph{does nonlinear signal extraction add value over the
regularized linear map, once costs are paid?} Phases 4--8 traded a single standardized
deviation through a linear threshold policy and found a small, cost-dominated edge. If the
relationship between the spread's level and its subsequent reversion is nonlinear, a
flexible model should recover value the linear rule leaves on the table. This phase tests
that, and reports the negative.

\subsection{The forecasting task}

The target is the realized $h$-bar ($h=5$) market-neutral return of fading a rich causal
rolling-OLS-120 spread:
\begin{equation}
  R_t^{h} \;=\; -\,\frac{S_{t+h}-S_t}{g_t},
  \qquad
  S_t = \log P^{T}_t - \hat\alpha_t - \sum_j \hat\beta_{j,t}\log P^{j}_t,
  \qquad
  g_t = 1+\sum_j\lvert\hat\beta_{j,t}\rvert ,
  \label{eq:nonlin_target}
\end{equation}
with $S_t$ the causal (one-bar-delayed) spread and $g_t$ the grossing factor
\eqref{eq:grossing} that converts a spread change into a return on \$1 gross. Predicting
$R_t^h>0$ means ``the rich spread will revert'', which is exactly the trade the threshold
book takes --- so the forecasting task and the trading task are the same task, and the
comparison is meaningful.

\paragraph{Features (11, all causal).} Current $z$; spread level; spread changes over 1
and 5 bars; 5-bar and 20-bar rolling volatility of the spread; target return and its first
two lags; basket return and its first two lags. Every feature uses bars $\le t$ only. The
set is deliberately small and low-capacity: with daily data and a weak signal,
high-dimensional inputs inflate variance without revealing structure
(Section~\ref{sec:nonlinear:trees}) --- a design decision made from
Chapter~\ref{ch:nonlinear}'s theory rather than from experimentation on the test set.

\paragraph{Estimators.} \code{linear}: Ridge on standardized features (the
Chapter~\ref{ch:regularization} baseline). \code{gbm}:
\code{HistGradientBoostingRegressor}. \code{mlp}: shallow \code{MLPRegressor} (run on
\code{ADBE} only, for compute reasons --- itself a finding: the nonlinear path is
$\approx100\times$ more expensive per fit). A \code{persistence} reference (the IC of
$z_t$ itself) is computed alongside, so ``how much does the model add beyond simply
trading the current mispricing?'' is directly measurable.

\paragraph{No-lookahead discipline.} Forecasters are fit on a causal \emph{expanding}
window: at each quarterly (63-bar) boundary the estimator is retrained on every valid
training bar from warm-up onward, then predicts only the strictly subsequent held-out
bars. Feature standardisation uses training-window mean and standard deviation applied to
the held-out bars --- the classic leakage point, explicitly avoided. Reported metrics are
out-of-sample on 2022+ ($n\approx1{,}164$--$1{,}165$ per basket).

\subsection{Results: forecast quality}

\input{tables/nonlinear.tex}

\begin{figure}[H]\centering
\includegraphics[width=0.92\linewidth]{figures/nonlinear_ic.png}
\caption{Out-of-sample information coefficient by basket and estimator (2022+). The linear
Ridge is highest on every basket; the GBM is second on every basket; the MLP is between
them and still below linear.}
\label{fig:nonlinear_ic}
\end{figure}

Measured (Table~\ref{tab:nonlinear_forecast}):

\begin{center}
\renewcommand{\arraystretch}{1.2}
\begin{tabular}{llrrr}
\toprule
Basket & Model & OOS IC & OOS $R^2$ & $n$ \\
\midrule
\code{ADBE} & linear & $\mathbf{0.341}$ & $\mathbf{+0.088}$ & 1{,}164 \\
\code{ADBE} & gbm & $0.255$ & $-0.051$ & 1{,}164 \\
\code{ADBE} & mlp & $0.267$ & $-0.049$ & 1{,}164 \\
\code{AVGO} & linear & $\mathbf{0.423}$ & $\mathbf{+0.173}$ & 1{,}165 \\
\code{AVGO} & gbm & $0.366$ & $+0.119$ & 1{,}165 \\
\code{GS}   & linear & $\mathbf{0.300}$ & $\mathbf{+0.078}$ & 1{,}164 \\
\code{GS}   & gbm & $0.217$ & $-0.052$ & 1{,}164 \\
\code{CAT}  & linear & $\mathbf{0.243}$ & $\mathbf{+0.041}$ & 1{,}164 \\
\code{CAT}  & gbm & $0.173$ & $-0.058$ & 1{,}164 \\
\bottomrule
\end{tabular}
\end{center}

\subsection{Results: is the linear equation misspecified?}

\input{tables/nonlinear_reset.tex}

\begin{figure}[H]\centering
\includegraphics[width=0.85\linewidth]{figures/nonlinear_reset.png}
\caption{Ramsey RESET $F$-test on the linear forecasting equation, fitted on the training
sample only. Three of four baskets reject linearity at 5\%; the \code{ADBE} basket used
throughout Phases 4--8 does not.}
\label{fig:nonlinear_reset}
\end{figure}

Measured (Table~\ref{tab:nonlinear_reset}, $n_{\text{train}}=2{,}881$ per basket):
\code{ADBE} $p=0.376$ (no rejection), \code{AVGO} $p<10^{-3}$, \code{GS} $p=0.0023$,
\code{CAT} $p<10^{-3}$ (all reject).

\subsection{Results: the learning curve, i.e.\ \emph{why}}

\input{tables/nonlinear_learn.tex}

\begin{figure}[H]\centering
\includegraphics[width=0.85\linewidth]{figures/nonlinear_learning.png}
\caption{GBM learning curve: out-of-sample $R^2$ on a fixed 2019--2020 validation window
as the causal training sample grows. The ensemble is deep inside its variance-dominated
regime until $\approx2{,}400$--$2{,}800$ training bars --- which is about as much data as
this universe provides.}
\label{fig:nonlinear_learning}
\end{figure}

Measured (Table~\ref{tab:nonlinear_learning}): OOS $R^2$ runs $-0.446$ at 400 training
bars, $-0.304$ at 800, $-0.228$ at 1{,}200, $-0.170$ at 1{,}600, $-0.132$ at 2{,}000,
$-0.031$ at 2{,}400, and only at 2{,}800 bars does it turn positive ($+0.432$ on that
particular window). With $\approx2{,}881$ training bars available in total, the flexible
model reaches competence at the very end of the available data and never gets to
demonstrate superiority. This is the textbook small-signal, low-$N$ regime of
Section~\ref{sec:nonlinear:trees}: boosting's bias is near zero and all its error is
variance, and variance needs data.

\subsection{Results: does any of it survive contact with costs?}

Forecast IC is a statistical summary; what matters is the traded result. So the
Phase~\ref{ch:phase4} $z$-threshold book on the \code{ADBE} basket is given an entry
\emph{gate}: a position opens only when the threshold condition holds \emph{and} the
forecaster's predicted $R_t^h$ points in the same reversion direction.

Table~\ref{tab:nonlinear_trade} (generated from
\code{results/tables/nonlinear\_trading.csv}) gives the full comparison; the headline
rows are: baseline $z$-threshold book $+0.050$ net Sharpe over 198 trades (full sample);
linear gate $+0.037$ over 98 trades; GBM gate $+0.003$ over 95; MLP gate $-0.018$ over 92.
On the 2022+ test period the baseline is $-0.302$ over 48 trades and \emph{every} gated
variant is $-0.346$ over 23 trades.

The verdict is negative and unambiguous: \textbf{gating never raises net Sharpe}, on
either the full sample or the test period. The linear gate approximately halves the number
of trades (198$\to$98) without improving net Sharpe ($0.050\to0.037$); the GBM and MLP
gates, whose marginal predictions are less accurate, degrade it further ($0.003$ and
$-0.018$). Filtering entries by a nonlinear reversion forecast therefore does not recover
the value the statistical tests hint at --- once transaction costs are paid.

\paragraph{Why halving the trades did not help.} \eqref{eq:costdrag} says the drag is
$c\cdot\mathrm{TO}$, so halving trades should halve the drag and add $\approx1.3\%$/year.
It did not, because the gate removes trades \emph{proportionally}: the filtered-out trades
were, on average, as profitable as the kept ones. That is the operational meaning of
``IC $=0.25$ but no incremental value over $z_t$ itself'' --- the nonlinear forecast is
correlated with the current mispricing (which is the first feature it is given!) and adds
almost no independent ranking information. The persistence reference in
\code{nonlinear\_summary} quantifies exactly this.

\subsection{Findings}

\begin{enumerate}[label=\textbf{N\arabic*.}]
  \item \textbf{Linear wins the forecasting race on every basket.} Highest OOS IC
    ($0.243$--$0.423$) and the only clearly positive OOS $R^2$ ($+0.041$ to $+0.173$). The
    GBM is second everywhere; the MLP is comparable to the GBM and below linear. Whatever
    functional nonlinearity these models can capture in-sample does not generalise at
    daily frequency on $\approx3{,}000$ bars.
  \item \textbf{The linear equation is not strongly misspecified where it matters.} RESET
    rejects linearity on three of four baskets but not on \code{ADBE} ($p=0.376$) --- the
    basket used throughout Phases 4--8. And where linearity \emph{is} rejected (AVGO, GS,
    CAT), the flexible models still fail to convert the misspecification into OOS $R^2$
    (GS: linear $+0.078$ vs GBM $-0.052$). Detecting nonlinearity and exploiting it are
    different problems; the second needs far more data than the first.
  \item \textbf{The learning curve explains the result without reference to the market.}
    The GBM's OOS $R^2$ is negative until $\approx2{,}400$--$2{,}800$ training bars, which
    is essentially the whole dataset. The failure is a sample-size failure, not evidence
    that markets are linear.
  \item \textbf{Nothing survives costs.} No gate improves net Sharpe on any scope, and the
    honest gate in \code{environment/check.sh} \emph{fails the build} if a future edit ever
    lets one: the check reads \code{nonlinear\_trading.csv} and asserts that the
    linear-gate net Sharpe does not exceed the baseline's.
\end{enumerate}

\subsection{Honest conclusion, and what would change it}

\textbf{Conclusion.} At daily frequency, on $\approx3{,}000$--$4{,}000$ observations per
basket, with a signal-to-noise ratio of a few percent, a regularized linear map dominates
gradient boosting and a shallow network both as a forecaster and as a trading gate. The
result is predicted in advance by the bias--variance decomposition
\eqref{eq:biasvariance}--\eqref{eq:mse:decomp} and confirmed by the learning curve; it is
not a surprise and it is not a comment on nonlinear models in general.

\textbf{What would change it} (Section~\ref{sec:nonlinear:when}): pooling all 195 names'
baskets into one supervised problem with basket fixed effects --- raising the effective
sample by two orders of magnitude while keeping per-basket economics; higher-frequency data
where SNR is better and $T$ per year is 250$\times$ larger; classification targets with
genuine asymmetry (``will this relationship break?'') rather than a regression on a weak
continuous return; and features with real nonlinear structure (liquidity, borrow cost,
option-implied skew). The first is the cheapest and is listed first in
Chapter~\ref{ch:conclusion}.

\subsection{Exit gate --- status}

\begin{itemize}
  \item[\checkmark] nonlinear versus linear OOS $R^2$, IC and net Sharpe compared with an
    honest conclusion;
  \item[\checkmark] learning curves and RESET plotted from generated CSVs;
  \item[\checkmark] pipeline unit tests pass (\code{scripts/test\_nonlinear.py});
  \item[\checkmark] the build \emph{fails} if the nonlinear gate ever outclaims the
    baseline net Sharpe (honesty gate in \code{check.sh}).
\end{itemize}
